\chapter*{Introduction}
\label{ch:intro}

% Introduzione generale: Presenta il tema, il problema di ricerca e gli obiettivi del lavoro.
The rapid evolution of Vision-Language Models raises the need to determine whether their strong performance reflects genuine multimodal understanding 
or the exploitation of superficial correlations and linguistic biases. In this contex \textit{MAIA: Multimodal AI Assessment} by Testa et al.(2025)\cite{testa2025maia,testa2025allinone} was introduced as a native Italian 
benchmark for evaluating video-based reasoning in Vision-Language Models.
Testa et al. investigate the extent to which VLMs coherently integrate visual and linguistic information. More specifically, it assesses whether model predictions 
are grounded in visual evidence rather than primarily driven by distributional regularities learned by the language component. This distinction is crucial because 
a model may answer a question correctly while failing to recognise the same information consistently when it is presented in a discriminative format. In such cases,
the prediction may result from \underline{linguistic plausibility} rather than genuine visual understanding.
To address this issue, MAIA employs two aligned tasks: Visual Statement Verification and Open-Ended Visual Question Answering. 
Through these tasks, Testa et al. evaluate not only models accuracy, but also the robustness and consistency of models behaviour across different evaluation formats.
The benchmark further defines twelve fine-grained reasoning categories to identify which competencies are used by the models and in which areas linguistic shortcuts, 
unimodal collapse, or hallucinations emerge.
A limitation of the original framework is that the dataset was designed and evaluated primarily on the basis of visual information. 
Real-world videos, however, are inherently audiovisual: relevant meaning may also be conveyed through dialogue, environmental sounds, 
music, and other acoustic events. Recent multimodal and omnimodal architectures can process visual, auditory, and textual inputs, but the availability of 
multiple modalities does not guarantee that they will be integrated effectively. Models may over-rely on a single channel or be negatively affected by redundant 
or weakly relevant information.
Starting from this problem, the study Deodati,Kaif et al.(2026): "\textit{Lost in Transcription: Investigating the Role of Audio Information for Video-based Visual
Statement Verification}"\cite{deodati2026lost} extends the original MAIA research framework to a more complex analysis of the interaction between visual, linguistic, 
and auditory information, while preserving its competence-oriented perspective.
The study considers two complementary representations of audio: raw signal and automatic transcriptions provided as textual input. 
By evaluating these representations both independently and in combination with video information, the analysis investigates whether audio contributes 
meaningfully to task performance and whether one representation has a stronger influence on model predictions.
This leads to two research questions:

\textbf{RQ1}: \textit{To what extent does raw audio affect video-based Visual Statement Verification when it is jointly processed with visual and textual information?}

\textbf{RQ2}: \textit{How effectively do multimodal models integrate visual, textual, and auditory information when performing video-based reasoning and understanding tasks?}

The experiments show that raw audio has a limited and strongly context-dependent effect. It is most useful when visual information is unavailable, whereas, when added to an already
informative visual input, it may provide little benefit or introduce noise. The results also show that transcriptions can produce greater interference than raw audio, suggesting 
that multimodal architectures do not necessarily achieve more effective cross-modal understanding simply because additional modalities are available.
These findings must nevertheless be interpreted in relation to the nature of the dataset. The one hundred original MAIA videos were selected and annotated according to their 
visual content, and the questions were formulated without considering the audio track. Consequently, the benchmark does not systematically require auditory information to solve its items.
My thesis project begins from this limitation. The following chapters investigate whether models can integrate multiple input modalities more coherently when the dataset 
is specifically designed so that auditory information is at least partially relevant. The central hypothesis is that, under these conditions, models may use complementary 
signals across modalities to reconstruct missing information and develop a more genuinely multimodal understanding of video content, producing answers grounded in the available 
evidence rather than in linguistic plausibility or chance.
